\vskip4ex
\wysubtitle{logos.pdf}
{\makebox[.5\textwidth][l]{出题：\textsf{\textbf{Math}Hub}}\makebox[.5\textwidth][l]{审题：未央学社}}
{https://qm.qq.com/q/OlkSiy9W04}{HDU 数学营}
{http://weixin.qq.com/r/mhKYgFTEMQlOrRDj90eI}{未央学社公众号}
\begin{flushleft}
    \kaishu\qquad 本资料仅作为模拟练习之用，目的是为了帮助大家更有效地复习，并减轻对考试的担忧。请正确的对待此资料，其旨在辅助复习，而非预示具体的考试内容。我们鼓励同学们认真复习，大学学习主打理解，而非刷题，期望大家在期中考试中取得优异成绩。
\end{flushleft}

\section{选择题}
\begin{choice}{A}
    设$y=\tan ^2x\sin \frac{1}{x}$，则当$x\to 0$时，$y$是
    \begin{tasks}(4)
        \task 无穷小量
        \task 无穷大量
        \task 有界但非无穷小量
        \task 无界但非无穷大量
    \end{tasks}
\end{choice}
\begin{solution}*
    本题考察了无穷小性质，无穷小乘以有界量为无穷小

    \textbf{解：}$x\to 0$时，$\tan ^2x\sin \frac{1}{x}\sim x^2\sin \frac{1}{x}\to 0$.\sokka{A}
\end{solution}

\begin{choice}{D}
    设函数$f(u)$可导，$y=f\ab(x^2)$当自变量$x$在$x=-1$处取得增量$\Delta x=-0.1$时，相应的函数的增量$\Delta y$的线性主部为0.1，则$f'(1) =$
    \begin{tasks}(4)
        \task $-1$
        \task $0.1$
        \task $1$
        \task $0.5$
    \end{tasks}
\end{choice}
\begin{solution}*
    本题考察了微分的概念
    
    \textbf{解：}$\d y=\ab[f\ab(x^2)] '\d x=f'\ab(x^2) \cdot 2x\d x=2xf'\ab(x^2) \d x$，即
    \[0.1=2\cdot (-1) \cdot f'\ab[(-1) ^2 ] \cdot (-0.1) \Longrightarrow 1=-2f'(1) \cdot (-1) \Longrightarrow 2f'(1) =1\Longrightarrow f'(1) =\frac{1}{2}\]
    \sokka{D}
\end{solution}

\begin{choice}{D}
    下列各式中正确的是
    \begin{tasks}(4)
        \task $\lim_{x\to +\infty} \frac{\sin x}{x}=1$
        \task $\lim_{x\to +\infty} \frac{\ln(1+x)}{x}=1$
        \task $\lim_{x\to 1} \frac{\arctan x}{x}=1$
        \task $\lim_{x\to 0} \frac{\e^x-1}{x}=1$
    \end{tasks}
\end{choice}
\begin{solution}*
    \begin{tasks}(2)
        \task\color{solutioncolor} $\abs{\sin x}\leq 1$，于是$\lim_{x\to +\infty} \frac{\sin x}{x}=0$.
        \task\color{solutioncolor} $\lim_{x\to +\infty} \frac{\ln(1+x)}{x}=0$，$\ln(1+x)$的速度远慢于$x$.
        \task\color{solutioncolor} $\lim_{x\to 1} \frac{\arctan x}{x}=\frac{\arctan 1}{1}=\frac{\frac{\pi}{4}}{1}=\frac{\pi}{4}$.
        \task\color{solutioncolor} $\lim_{x\to 0} \frac{\e^x-1}{x}=\lim_{x\to 0} \frac{x}{x}=1$.
    \end{tasks}
\sokka{D}
\end{solution}

\begin{choice}{B}
   设$f'(0) =2$，则$\lim_{x\to 0} \frac{f(3\sin x)-f(2\arctan x)}{x}=$
   \begin{tasks}(4)
    \task $\frac{1}{2}$
    \task 2
    \task $\frac{1}{4}$
    \task 4
    \end{tasks}
\end{choice}
\begin{solution}*
    本题考察了导数的定义，具体计算如下
    \[
    \begin{aligned}
        L &= \lim_{x\to 0} \frac{f(3\sin x)-f(2\arctan x)}{x}= \lim_{x\to 0} \frac{f(3\sin x)-f(0) +f(0) -f(2\arctan x)}{x} \\
          &= \lim_{x\to 0} \frac{f(3\sin x)-f(0) -\ab[f(2\arctan x) -f(0) ]}{x}= \lim_{x\to 0} \frac{f(3\sin x)-f(0)}{x}-\lim_{x\to 0} \frac{f(2\arctan x) -f(0)}{x} \\
          &= \lim_{x\to 0} \frac{f(0+3\sin x)-f(0)}{3\sin x}\cdot \frac{3\sin x}{x}-\lim_{x\to 0} \frac{f(0+2\arctan x)-f(0)}{2\arctan x}\cdot \frac{2\arctan x}{x} \\
          &= \lim_{x\to 0} \frac{f(0+3\sin x)-f(0)}{3\sin x}\cdot \lim_{x\to 0} \frac{3\sin x}{x}-\lim_{x\to 0} \frac{f(0+2\arctan x)-f(0)}{2\arctan x}\cdot \lim_{x\to 0} \frac{2\arctan x}{x} \\
          &= f'(0) \lim_{x\to 0} \frac{3x}{x}-f'(0) \lim_{x\to 0} \frac{2x}{x}= 3f'(0) -2f'(0)= f'(0) = 2
    \end{aligned}
    \]
    \sokka{B}
\end{solution}

\begin{choice}{D}
    设$f'(x_0) =f''(x_0) =0$，$f'''(x_0) >0$，则下列选项正确的是
    \begin{tasks}(2)
        \task $f'(x_0) $是$f'(x)$的极大值
        \task $f(x_0)$是$f(x)$的极大值
        \task $f(x_0)$是$f(x)$的极小值
        \task $\ab(x_0,f(x_0))$是曲线$y=f(x)$的拐点
    \end{tasks}
\end{choice}
\begin{solution}*
    本题考察了微分的概念本题考察了极限保号性，导数的定义、二阶导的正负性和极值、拐点的关系
    
    \textbf{解：}由导数定义：$f'''(x_0) =\lim_{x\to x_0} \frac{f''(x) -f''(x_0)}{x-x_0}=\lim_{x\to x_0} \frac{f''(x)}{x-x_0}>0$.
    
    由极限的保号性可知，在去心邻域$\mathring{U}(x_0,\delta)$中$\left\{\begin{aligned}
        &\text{当} x\in (x_0,\delta)\text{ 时，}x-x_0>0,\ f''(x) >0\\
        &\text{当} x\in (-\delta,x_0)\text{ 时，}x-x_0<0,\ f''(x) <0
    \end{aligned}\right.$.

    于是，$\ab(x_0,f(x_0))$是曲线$y=f(x)$的拐点.
    \sokka{D}
\end{solution}

\begin{choice}{C}
    设函数$g(x)$可导，$h(x) =\e^{1+g(x)}$，其中$h'(1) =1,g'(1) =2$，则$g(1) =$ 

    \begin{tasks}(4)
    \task $\ln 3-1$
    \task $-\ln 3-1$
    \task $-\ln 2-1$
    \task $\ln 2-1$
    \end{tasks}
\end{choice}
\begin{solution}*
    本题考察的是用求导公式和求导法则求抽象复合函数的导数

    \textbf{解：}由题意可知：
    $h'(x) =\e^{1+g(x)}\cdot\ab[1+g(x) ] '=\e^{1+g(x)}\cdot\ab[0+g'(x) ] =g'(x) \e^{1+g(x)}$.

    将$h'(1) =1,g'(1) =2$代入其中有
    \[
    \begin{aligned}
        h'(1) = g'(1) \e^{1+g(1)}&\Longrightarrow 1 = 2\e^{1+g(1)}\Longrightarrow \e^{1+g(1)} = \frac{1}{2}\Longrightarrow \ln\e^{1+g(1)} = \ln\ab(\frac{1}{2})\\
        &\Longrightarrow 1+g(1) = \ln 2^{-1}\Longrightarrow g(1) = -\ln 2 - 1
    \end{aligned}    
    \]
    \sokka{C}
\end{solution}

\begin{choice}{C}
    曲线$y_{1}=ax^3+1$与$y_{2}=\e^{x}$在$x=1$处斜率相等，则$a=$
    \begin{tasks}(4)
        \task $\e $
        \task $\frac{\e }{2}$
        \task $\frac{\e }{3}$
        \task $\frac{\e }{4}$
    \end{tasks}
\end{choice}
\begin{solution}*
    本题考察了切线斜率与导数的关系.
    
    \textbf{解：}若两曲线在$x=1$处斜率相等，则在该点处的导数相同
    \[\begin{aligned}
        \begin{cases}
            y_{1}'=3ax^2+0=3ax^2\\
            y_{2}'=\e^x\\
        \end{cases}\Longrightarrow
        \begin{cases}
            y_{1}'(1) =3a\cdot 1^2=3a\\
            y_{2}'(1) =\e^1=\e \\
        \end{cases}\Longrightarrow y_{1}'(1) =y_{2}'(1)
        \Longrightarrow 3a=\e \Longrightarrow a=\frac{\e }{3}.
    \end{aligned}\]
    \sokka{C}
\end{solution}

\begin{choice}{A}
    设$f(x) =\lim_{n\to \infty} \frac{(n-1)x}{nx^2+1}$，则$f(x) $的间断点为$x=$ 
    \begin{tasks}(4)
        \task 0
        \task 1
        \task 2
        \task 3
    \end{tasks}
\end{choice}
\begin{solution}*
    本题首先是利用极限求出函数表达式，再求得间断点

    \textbf{解：}
    \[\underset{n\to \infty}{\lim}\frac{nx-x}{nx^2+1}=\underset{n\to \infty}{\lim}\frac{x-\frac{x}{n}}{x^2+\frac{1}{n}}=\underset{n\to \infty}{\lim}\frac{1-\frac{1}{n}}{x+x\frac{1}{n}}=\frac{1}{x}\]
    $\lim_{x\to 0} \frac{1}{x} = \infty$，因此$ f(x)$ 的间断点为$x=0$.于是：$\ab(x_0,f(x_0))$是曲线$y=f(x)$的拐点.\sokka{A}
\end{solution}

\section{填空题}

\begin{problem}
    求极限$L=\lim_{x\to 0} \frac{x\ln (1+x)}{1-\cos x}=$\ans{ 2 }.
\end{problem}

\begin{draft}
    \vspace{3em}
\end{draft}
\begin{solution}*
    本题利用了等价无穷小替代求极限

    \textbf{解：}
    \[L=\lim_{x\to 0} \frac{x\ln (1+x)}{1-\cos x}=\lim_{x\to 0} \frac{x\cdot x}{\frac{1}{2}x^2}=2\lim_{x\to 0} \frac{x^2}{x^2}=2\]
\end{solution}

\begin{problem}
    求参数方程的导数：$\begin{cases}
        x=\cos t\\
        y=\sin t+t\\
    \end{cases}$求解结果=\ans{  $-\cot t-\csc t$  }.
\end{problem}
\begin{draft}
    \vspace{3em}
\end{draft}
\begin{solution}*
    本题考察了参数方程的导数：$\frac{\d y}{\d x}=\frac{\d y/\d t}{\d x/\d t}$.
    
    \textbf{解：}$\frac{\d x}{\d t}=-\sin t,\ \frac{\d y}{\d t}=\cos t+1$.于是有

    \[y'=\frac{\d y}{\d x}=\frac{\d y/\d t}{\d x/\d t}=\frac{\cos t+1}{-\sin t}=-\frac{\cos t}{\sin t}-\frac{1}{\sin t}=-\cot t-\csc t\]
\end{solution}


\begin{problem}
    求曲线$y=x\e^{-x}$的拐点及在拐点处的曲率分别为\ans{  $\ab(2,2\e^{-2})$，$\kappa = 0$  }.
\end{problem}
\begin{draft}
    \vspace{3em}
\end{draft}

\begin{solution}*
    本题考察了拐点的概念和曲率公式
    
    \textbf{解：}
    \[
    \begin{aligned}
        y'&=\e^{-x}+x\ab(-\e^{-x} )=\e^{-x}-x\e^{-x}=\e^{-x}(1-x)\\
        y''&=-\e^{-x}(1-x) +\e^{-x}\cdot(-1)=\e^{-x}(x-1)-\e^{-x}=\e^{-x}(x-2)\\
    \end{aligned}    
    \]

    $y''=0\Longrightarrow \e^{-x}(x-2) =0\Longrightarrow x=2$
\begin{itemize}
    \item 当$x>2$时，$y''>0$；当$x<2$时，$y''<0$. 故$x=2$为拐点的横坐标.
    \item 当$x=2$时，$y=2\e^{-2}$. 故拐点为$( 2,2\e^{-2})$.
    \item 拐点处的曲率
    \[\kappa =\frac{\abs{y''}}{\ab[1+(y') ^2 ] ^{\frac{3}{2}}}=\frac{0}{\ab[1+(y') ^2 ] ^{\frac{3}{2}}}=0\]
\end{itemize}
\end{solution}

\section{解答题}
\begin{problem}
    求极限$L=\lim_{x\to 3} \frac{x-3}{x^2-9}$. 
\end{problem}
\begin{draft}
    \vspace{5em}
\end{draft}
\begin{solution}*
    分母利用平方差公式，再化简即可直接求解该极限

    \textbf{解：}$L=\lim_{x\to 3} \frac{x-3}{x^2-9}=\lim_{x\to 3} \frac{x-3}{(x-3)(x+3)}=\lim_{x\to 3} \frac{1}{x+3}=\frac{1}{3+3}=\frac{1}{6}$.
\end{solution}



\begin{problem}
    求极限$L=\lim_{n\to \infty}\ab[\frac{1}{(n+1)^2}+\frac{1}{(n+2) ^2}+\cdots +\frac{1}{(2n)^2} ]$. 
\end{problem}
\begin{draft}
    \vspace{5em}
\end{draft}
\begin{solution}*
    本极限利用夹逼准则求解

    \textbf{解：}首先对$\frac{1}{(n+1)^2}+\frac{1}{(n+2) ^2}+\cdots +\frac{1}{(2n)^2}$进行适当放缩：

    \[
    \begin{aligned}
        \frac{1}{(n+1)^2}+\frac{1}{(n+2) ^2}+\cdots +\frac{1}{(2n)^2}&\geqslant \frac{1}{(2n)^2}+\frac{1}{(2n)^2}+\cdots +\frac{1}{(2n)^2}=\frac{n}{(2n)^2}=\frac{n}{4n^2}=\frac{1}{4n}\\
        \frac{1}{(n+1)^2}+\frac{1}{(n+2) ^2}+\cdots +\frac{1}{(2n)^2}&\leq \frac{1}{(n+1)^2}+\frac{1}{(n+1)^2}+\cdots +\frac{1}{(n+1)^2}\\
        &=\frac{n}{(n+1)^2}=\frac{n}{n^2+2n+1}=\frac{1}{n+2+\frac{1}{n}}
    \end{aligned}
    \]

    而$\lim_{n\to \infty}\ab(\frac{1}{4n} )=0$且$\lim_{n\to \infty}\ab(\frac{1}{n+2+\frac{1}{n}} )=0$，由夹逼定理得$L=\lim_{n\to \infty}\ab[\frac{1}{(n+1)^2}+\frac{1}{(n+2) ^2}+\cdots +\frac{1}{(2n)^2} ] =0$.
\end{solution}


\begin{problem}
    求极限$L=\lim_{n\to \infty} \frac{n}{\ln n}\cdot\ab(\sqrt[n]{n}-1)$. 
\end{problem}
\begin{draft}
    \vspace{5em}
\end{draft}
\begin{solution}*
    本极限利用了幂指代换与等价替代求解；且$\exp f(x) =\e^{f(x)}$

    \textbf{解：}
    \[\begin{aligned}
        L&=\lim_{n\to \infty} \frac{n}{\ln n}\cdot\ab(\sqrt[n]{n}-1 )=\lim_{n\to \infty} \frac{n}{\ln n}\cdot\ab(n^{\frac{1}{n}}-1)=\lim_{n\to \infty} \frac{n}{\ln n}\cdot\ab(\e^{\frac{1}{n}\ln n}-1)\\
        &=\lim_{n\to \infty} \frac{n}{\ln n}\cdot \frac{1}{n}\ln n=\lim_{n\to \infty} \frac{n}{\ln n}\cdot \frac{\ln n}{n}=1
    \end{aligned}\]
\end{solution}

\begin{problem}
    设$y=f(\ln x)\e^{f(x)}$，其中$f$可微，则$\d y$为？
\end{problem}
\begin{draft}
    \vspace{5em}
\end{draft}
\begin{solution}*
    本题考察了复合函数求导(求微分)的问题

    \textbf{解：}
    \[\begin{aligned}
        y'&=\ab[f(\ln x)] '\cdot \e^{f(x)}+f(\ln x)\ab[\e^{f(x)} ]'=f'(\ln x)\cdot \frac{1}{x}\cdot \e^{f(x)}+f(\ln x)\cdot \e^{f(x)}\cdot f'(x)\\
        &=\e^{f(x)}\ab[\frac{f'(\ln x)}{x}+f(\ln x)f'(x) ]\Longrightarrow\d y=y'\d x=\e^{f(x)}\ab[\frac{f'(\ln x)}{x}+f(\ln x)f'(x) ] \d x
    \end{aligned}\]
\end{solution}

\begin{problem}
    设函数$y=y(x)$ 由方程$\ln\ab(x^2+y )=x^3y+\sin x$确定，则$\frac{\d y}{\d x}\Bigg|_{x=0}$为？
\end{problem}
\begin{draft}
    \vspace{5em}
\end{draft}
\begin{solution}*
    本题考察了隐函数求导，将代入$x=0,y=1$即可.

    \textbf{解：}当$x=0$时，代入方程中有
    \[\ln\ab(0^2+y )=0^3\cdot y+\sin 0\Longrightarrow \ln y=0\Longrightarrow y=1\]
    对方程$\ln\ab(x^2+y )=x^3y+\sin x$两边同时求导有
    \[\frac{1}{x^2+y}\cdot\ab(2x+y' )=3x^2y+x^3y'+\cos x\]
    将$x=0,\ y=1$代入上述方程有
    \[\begin{aligned}
        &\frac{1}{0^2+1}\cdot\ab(2\cdot 0+y'\Big|_{x=0,y=1}^{} )=3\cdot 0^2\cdot 1+0^3\cdot y'+\cos 0\\
        &\Longrightarrow y'\Big|_{x=0,y=1}^{}=1\Longrightarrow \frac{\d y}{\d x}\Bigg|_{x=0}^{}=1
    \end{aligned}\]
\end{solution}


\begin{problem}
    $y=x^{x^x}$，求$y'$.
\end{problem}
\begin{draft}
    \vspace{5em}
\end{draft}
\begin{solution}*
    本题考察了对数求导法.

    \textbf{解：}对$y=x^{x^x}$两边同时取对数有
    \[\begin{aligned}
        \ln y=\ln x^{x^x}&\Longrightarrow \ln y=x^x\ln x\\
        &\Longrightarrow \ln\ln y=\ln\ab(x^x\ln x )=\ln\ab(x^x )+\ln\ln x\\
        &\Longrightarrow \ln\ln y=x\ln x+\ln\ln x
    \end{aligned}\]
    两边求导有
    \[\begin{aligned}
        \frac{1}{\ln y}\cdot \frac{1}{y}\cdot y'&=\ln x+x\cdot \frac{1}{x}+\frac{1}{\ln x}\cdot \frac{1}{x}=\ln x+1+\frac{1}{x\ln x}\\
        &\Longrightarrow y'=y\ln y\ab(\ln x+1+\frac{1}{x\ln x} )
    \end{aligned}\]
    即
    \[\begin{aligned}
        y'&=x^{x^x}\cdot x^x\ln x\ab(\ln x+1+\frac{1}{x\ln x})\\
        &=x^{x^x}\cdot x^{x-1}\ab(x\ln ^2x+x\ln x+1)\\
        &=x^{x^x+x-1}\ab(x\ln ^2x+x\ln x+1)
    \end{aligned}
    \]
\end{solution}

\begin{problem}
    求极限：$L=\lim_{x\to 2} \frac{\sqrt{5x-1}-\sqrt{2x+5}}{x^2-4}$.
\end{problem}
\begin{draft}
    \vspace{5em}
\end{draft}
\begin{solution}*
    本极限利用了分子有理化求解.

    \textbf{解：}
        \[\begin{aligned}
            L&=\lim_{x\to 2} \frac{\sqrt{5x-1}-\sqrt{2x+5}}{x^2-4}=\lim_{x\to 2} \frac{\ab(\sqrt{5x-1}-\sqrt{2x+5})\ab(\sqrt{5x-1}+\sqrt{2x+5})}{\ab(x^2-4)\ab(\sqrt{5x-1}+\sqrt{2x+5})}\\
            &=\lim_{x\to 2} \frac{\ab(\sqrt{5x-1} )^2-\ab(\sqrt{2x+5} )^2}{\ab(x^2-4)\ab(\sqrt{5\cdot 2-1}+\sqrt{2\cdot 2+5})}
            =\lim_{x\to 2} \frac{\ab(5x-1 )-\ab(2x+5)}{\ab(x^2-4)\ab(\sqrt{9}+\sqrt{9})}\\
            &=\lim_{x\to 2} \frac{3x-6}{\ab(x^2-4)\ab(3+3)}
            =\lim_{x\to 2} \frac{3(x-2)}{6(x-2) (x+2)}=\lim_{x\to 2} \frac{1}{2(x+2)}=\frac{1}{2\times(2+2)}=\frac{1}{8}
        \end{aligned}\]
\end{solution}

\begin{problem}
    设$f(x)=(x^3-1)^n\e^x$，求$f^{(n)}(1)$.
\end{problem}
\begin{draft}
    \vspace{5em}
\end{draft}
\begin{solution}*
本题利用莱布尼茨求导公式即可，因为求的是$x=1$处的$n$阶导数，只要保留$[(x^3-1)^n]^{(n)}$项即可.

\textbf{解答：}根据莱布尼茨n阶导数公式有：$[u(x)v(x)]^(n)=\sum_{k=0}^n\binom{n}{k}u^{(n-k)}v^k$.

设$u(x)=(x^3-1)^n,\ v(x)=\e^x$，则有
\[\begin{aligned}
    f^{(n)}(x)&=\sum_{k=0}^n\binom{n}{k}u^{(n-k)}v^k\\
    &=\binom{n}{0}[(x^3-1)^n]^{(n)}\cdot\e^x+\binom{n}{1}\ab[(x^3-1)^n]^{(n-1)}\cdot(e^x)'+\cdots+\binom{n}{n}\ab[(x^3-1)^n]\cdot(e^x)^{(n)}\\
    &=\ab[(x^3-1)^n]^{(n)}\bigg|_{x=1}\cdot\e^1+\underset{n\text{个}0}{\underbrace{0+0+\cdots+0}}=\ab[(x^3-1)^n]^{(n)}\Bigg|_{x=1}\cdot\e.
\end{aligned}\]
下面单独研究$[(x^3-1)^n]^{(n)}$：
\[
\begin{aligned}
    \ab[(x^3-1)^n]'&=n\cdot(x^3-1)^{n-1}\cdot 3x^2\\
    \ab[(x^3-1)^n]''&=n(n-1)\cdot(x^3-1)^{n-2}\cdot 3x^2\cdot 3x^2+n\cdot(x^3-1)^(n-1)\cdot 6x\\
    &=n(n-1)\cdot(x^3-1)^{n-2}\cdot 3^2\cdot x^4+\cdots\\
    \ab[(x^3-1)^n]'''&=n(n-1)(n-2)\cdot(x^3-1)^{n-3}\cdot 3^3\cdot x^6+\cdots\\
    &\ \vdots\text{\quad（注意到除了第一项之外，求$n$阶导数均会保留$x^3-1$，代入$x=1$均为$0$）}\\
    \ab[(x^3-1)^n]^{(n-1)}&=n\cdot(n-1)\cdot(n-2)\cdot\cdots\cdot2\cdot(x^3-1)\cdot e^{n-1}\cdot x^{2(n-1)}+\cdots=n!\cdot(x^3-1)\cdot 3^{n-1}\cdot x^{2(n-1)}+\cdots\\
    \ab[(x^3-1)^n]^{(n)}&=n!\cdot 3x^2\cdot 3^{n-1}\cdot x^{2(n-1)}+\cdots=n!\cdot 3^n\cdot x^{2n}+\cdots\xlongequal{x=1}3^nn!\cdot 1^{2n}+0+\cdots 0=3^nn!\\
\end{aligned}
\]
综上所述：$f^{(n)}(1)=3^nn!\e$.
\end{solution}

\begin{problem}
    设$y=\frac{\e^x+1}{\e^x-1}$，求
    \begin{enumerate}
        \item 函数的间断点并判断其类型.
        \item 该函数的水平渐近线和垂直渐近线.
    \end{enumerate}
\end{problem}
\begin{draft}
    \vspace{5em}
\end{draft}
\begin{solution}*
    本题考察了间断点和渐近线的相关知识点

\begin{enumerate}
    \item 当$x=0$时，函数无定义，因此$x=0$为间断点.
    \item 当$x=0$时，函数无定义，$\lim_{x\to 0} \frac{\e^x+1}{\e^x-1}=\infty$，因此$x=0$为垂直渐近线.
    \[\begin{aligned}
        \lim_{x\to +\infty} f(x)&=\lim_{x\to +\infty} \frac{\e^x+1}{\e^x-1}=\lim_{x\to +\infty} \frac{1+\frac{1}{\e^x}}{1-\frac{1}{\e^x}}=\frac{1+0}{1-0}=1\\
        \lim_{x\to -\infty} f(x)&=\lim_{x\to -\infty} \frac{\e^x+1}{\e^x-1}=\frac{0+1}{0-1}=-1
    \end{aligned}\]
    因此，$y=-1,y=1$为水平渐近线.
\end{enumerate}
\end{solution}

\begin{problem}
    讨论函数$f(x)=\begin{cases}
        \frac{x}{1+\e^{\frac{1}{x}}},&x\neq 0\\
        0,&x=0
    \end{cases}$在点$x=0$的连续性和可导性质.
\end{problem}
\begin{draft}
    \vspace{5em}
\end{draft}
\begin{solution}*连续性通过讨论左右极限和函数值；可导性讨论左右导数.
\[
\begin{cases}
    \text{左极限：}f(0-0)=\lim_{x\to 0^-}\frac{x}{1+\e^{\frac{1}{x}}}=\frac{0}{1+0}=0\\
    \text{函数值：}f(0)=0\\
    \text{右极限：}f(0+0)=\lim_{x\to 0^+}\frac{x}{1+\e^{\frac{1}{x}}}=\frac{0}{1+\infty}=0\\
\end{cases}    
\]
于是，$f(0-0)=f(0)=f(0+0)\Longrightarrow f(x)$在$x=0$处连续.

下面讨论$f(x)$在$x=0$处的可导性质,
\[\left\{
\begin{aligned}
    \text{左导数：}f_-'(0)&=\lim_{x\to 0^-}\frac{f(x)-f(0)}{x-0}=\lim_{x\to 0^-}\frac{\frac{x}{1+\e^{\frac{1}{x}}}-0}{x}=\lim_{x\to 0^-}\frac{x}{\ab(1+\e^{\frac{1}{x}})}\\
    &=\lim_{x\to 0^-}\frac{1}{1+\e^{\frac{1}{x}}}=\frac{1}{1+0}=1\\
    \text{右导数：}f_+'(0)&=\lim_{x\to 0^+}\frac{f(x)-f(0)}{x-0}=\lim_{x\to 0^+}\frac{\frac{x}{1+\e^{\frac{1}{x}}}-0}{x}=\lim_{x\to 0^+}\frac{x}{\ab(1+\e^{\frac{1}{x}})}\\
    &=\lim_{x\to 0^+}\frac{1}{1+\e^{\frac{1}{x}}}=\frac{1}{1+\infty}=0
\end{aligned}
\right.\]
于是：$f_-'(0)\neq f_+'(0)\Longrightarrow f(x)$在$x=0$处不可导.
\end{solution}

\begin{problem}
    设$0<x<a$，对任意的自然数$m,n$，试证：$x^m(a-x) ^n\leq \frac{m^mn^n}{(m+n) ^{m+n}}a^{m+n}$.
\end{problem}
\begin{draft}
    \vspace{5em}
\end{draft}
\begin{solution}*
    本题通过研究左边函数的单调性，确认出极值（最值），再利用其与不等式右边的数相比较

\columnseprule=0.5pt
\begin{multicols}{2}
    \textbf{证明}：设$f(x) =x^m(a-x) ^n$

    \[\begin{aligned}
        f'(x)&=mx^{m-1}(a-x) ^n-nx^m(a-x) ^{n-1}\\
        &=x^{m-1}(a-x) ^{n-1}\ab[m(a-x) -nx ]\\
        &=x^{m-1}(a-x) ^{n-1}\ab[ma-(m+n) x ]
    \end{aligned}\]

    令$f'(x) =0$，则$x^{m-1}(a-x) ^{n-1}\ab[ma-(m+n) x ] =0$.
    
    $x\in (0,a)$时，$x^{m-1}(a-x) ^{n-1}\ne 0$，故$ma-(m+n) x=0$，即$x=\frac{ma}{m+n}$.下面进行列表

\begin{center}
    \begin{tabular}{c|c|c|c}
        \toprule
        $x$&		$\ab(0,\frac{ma}{m+n})$&		$\frac{ma}
        {m+n}$&		$\ab(\frac{ma}{m+n},a)$\\
        \midrule
        $f'(x)$&		$+$&		$0$&		$-$\\
        $f(x)$&		$\nearrow$&	极大值&$\searrow$\\
        \bottomrule
    \end{tabular}
\end{center}
\end{multicols}
        
    因此$f(x)$在$x=\frac{ma}{m+n}$取得极大值：

    \[\begin{aligned}
        f\ab(\frac{ma}{m+n} )&=\ab(\frac{ma}{m+n} )^m\ab(a-\frac{ma}{m+n} )^n=\frac{m^ma^m}{(m+n) ^m}\ab[\frac{a(m+n)}{m+n}-\frac{ma}{m+n} ] ^n\\
        &=\frac{m^ma^m}{(m+n) ^m}\ab(\frac{na}{m+n} )^n=\frac{m^ma^m}{(m+n) ^m}\cdot \frac{n^na^n}{(m+n) ^n}=\frac{m^mn^n}{(m+n) ^{m+n}}a^{m+n}
    \end{aligned}\]

    又因为$f(0) =f(a) =0$，故$f\ab(\frac{ma}{m+n})$也为$x\in (0,a)$时，$f(x)$的最大值. 因此，在$x\in (0,a)$时，$f(x) \leq \frac{m^mn^n}{(m+n) ^{m+n}}a^{m+n}$，即不等式$x^m(a-x) ^n\leq \frac{m^mn^n}{(m+n) ^{m+n}}a^{m+n}$成立.
\end{solution}

\begin{problem}
    设$f(x)$在$[0,+\infty)$上可导，且$0\leq f(x) \leq \frac{x}{1+x^2}$.证明：存在$\xi \in (0,+\infty)$使得$f'(\xi) =\frac{1-\xi ^2}{\ab(1+\xi ^2 )^2}$.
\end{problem}
\begin{draft}
\begin{minipage}{0.81\textwidth}
\begin{flushright}\color{black}
        扫码观看讲解\href{https://b23.tv/W2KQlsI}{【未央学社】2023杭电高等数学上期中模拟卷一讲解}

        答案于微信公众号~\href{http://weixin.qq.com/r/mhKYgFTEMQlOrRDj90eI}{WY未央学社}~后台回复“\textbf{2023期中模拟卷}”获取
\end{flushright}
\end{minipage}
\begin{minipage}{0.15\textwidth}
        \rightline{\qrcode{https://b23.tv/W2KQlsI}\quad}
\end{minipage}
\end{draft}
\begin{solution}*
    构造函数后利用罗尔定理证明即可

    \textbf{证明：}设
    $F(x) =f(x) -\frac{x}{1+x^2}$，$F(x)$在$[0,+\infty)$连续且可导. 由于$0\leq f(x) \leq \frac{x}{1+x^2}$可得
\[f(0) =0,\ f(+\infty) =\lim_{x\to +\infty} f(x) \leq \lim_{x\to +\infty} \frac{x}{1+x^2}=\lim_{x\to +\infty} \frac{1}{1+\frac{1}{x}}=0\]
于是$F(0) =f(0) -0=0,\ F(+\infty) =\lim_{x\to +\infty}\ab[f(x) -\frac{x}{1+x^2} ] =f(+\infty) -\lim_{x\to +\infty} \frac{x}{1+x^2}=0$. 由于$F(0) =F(+\infty) =0$，由罗尔定理，存在$\xi \in (0,+\infty)$使得
\[f'(\xi) =0\Longrightarrow f(\xi) -\frac{\xi}{1+\xi ^2}=0\Longrightarrow f'(\xi) =\frac{1-\xi ^2}{\ab(1+\xi^2)^2}\]
\footnotesize
    \begin{spacing}{0.5}
        \textbf{构造函数：}
        令$\xi =x$，则有$f'(x) =\frac{1-x^2}{\ab(1+x^2)^2}\Longrightarrow f'(x) -\frac{1-x^2}{\ab(1+x^2)^2}=0$.注意到$f(x) =\frac{x}{1+x^2}$的导数为
        \[f'(x) =\frac{1\cdot \ab(1+x^2)-x\cdot 2x}{\ab(1+x^2)^2}=\frac{1+x^2-2x^2}{\ab(1+x^2)^2}=\frac{1-x^2}{\ab(1+x^2)^2}\Longrightarrow F(x) =f(x) -\frac{x}{1+x^2}\]
    \end{spacing}
\end{solution}